%\documentclass[twocolumn,superscriptaddress,aps,pra,showpacs]{revtex4}
\documentclass[prl,twocolumn,superscriptaddress,floatfix,showpacs,amsmath,longbibliography,nofootinbib,amssymb]{revtex4-1}
\usepackage[colorlinks=true,citecolor=blue,linkcolor=blue]{hyperref}
\usepackage[usenames]{color}
\usepackage{subfigure}
\usepackage{dcolumn}
\usepackage{mathrsfs}% Align table columns on decimal point
\usepackage{bm}% bold math
\usepackage{amsmath,amssymb,epsfig,float,graphics}
%\DeclareMathOperator{\atanh}{atanh}
\DeclareMathOperator{\diag}{diag}
\usepackage{graphicx}% Include figure files
\usepackage{dcolumn}% Align table columns on decimal point
\usepackage{bm}% bold math
\usepackage{hyperref}% add hypertext capabilities
\usepackage{xcolor}
\usepackage{amsmath, amssymb, amsfonts}
\usepackage{mathtools}
\usepackage{subfigure}
\usepackage{mathrsfs}
\newcommand{\col}[1]{\textcolor{red}{#1}}
\newcommand{\colb}[1]{\textcolor{blue}{#1}}
\newcommand{\colg}[1]{\textcolor{orange}{#1}}
\definecolor{purple1}{rgb}{128,0,128}
\newcommand{\colp}[1]{\textcolor{purple}{#1}}
\newcommand{\ket}[1]{\left|#1\right>}
\newcommand{\bra}[1]{\left<#1\right|}
\newcommand{\nn}{\nonumber\\}
\newcommand{\bea}{\begin{eqnarray}}
\newcommand{\ea}{\end{eqnarray}}
\newcommand{\ord}{{\cal O}}
\newcommand{\qq}{{\bf ???}}
\DeclareMathOperator{\sech}{sech}
%\newcommand{\bra}[1]{\langle#1|}
%\newcommand{\ket}[1]{|#1\rangle}
%\usepackage{graphicx}
%\usepackage{amsfonts}
%\usepackage{hyperref}
%\usepackage{mathrsfs}
%\usepackage{bbm}
%\usepackage{times,txfonts}
%\usepackage{isomath}
%\usepackage{units}
%\usepackage{nicefrac}


\begin{document}
\baselineskip=0.45 cm

\title{Verifying the upper bound on the speed of 
scrambling with the analogue Hawking radiation of trapped ions}

\author{Zehua Tian}
\email{tianzh@ustc.edu.cn}
\affiliation{Hefei National Laboratory for Physical Sciences at the Microscale and Department of Modern Physics, University of Science and Technology of China, Hefei 230026, China}
\affiliation{CAS Key Laboratory of Microscale Magnetic Resonance, University of Science and Technology of China, Hefei 230026, China}
\affiliation{Synergetic Innovation Center of Quantum Information and Quantum Physics, University of Science and Technology of China, Hefei 230026, China}
%\affiliation{Key Laboratory for Research in Galaxies and Cosmology, Chinese Academy of Science, 96 JinZhai Road, Hefei 230026, Anhui, China}




\author{Yiheng Lin}
\email{Yiheng@ustc.edu.cn}
\affiliation{Hefei National Laboratory for Physical Sciences at the Microscale and Department of Modern Physics, University of Science and Technology of China, Hefei 230026, China}
\affiliation{CAS Key Laboratory of Microscale Magnetic Resonance, University of Science and Technology of China, Hefei 230026, China}
\affiliation{Synergetic Innovation Center of Quantum Information and Quantum Physics, University of Science and Technology of China, Hefei 230026, China}


\author{Uwe R. Fischer}
\affiliation{Seoul National University, Department of Physics and Astronomy, Center for Theoretical Physics, Seoul 08826, Korea}

%\author{Andrzej Dragan}
%\affiliation{Institute of Theoretical Physics, University of Warsaw, Pasteura 5, 02-093 Warsaw, Poland}


\author{Jiangfeng Du}
\email{djf@ustc.edu.cn}
\affiliation{Hefei National Laboratory for Physical Sciences at the Microscale and Department of Modern Physics, University of Science and Technology of China, Hefei 230026, China}
\affiliation{CAS Key Laboratory of Microscale Magnetic Resonance, University of Science and Technology of China, Hefei 230026, China}
\affiliation{Synergetic Innovation Center of Quantum Information and Quantum Physics, University of Science and Technology of China, Hefei 230026, China}







\begin{abstract}
A general bound on the Lyapunov exponent of a %chaotic 
quantum system is given by 
$\lambda_L\leq2\pi\,T/\hbar$, where $T$  is the system temperature, as established by 
Maldacena, Shenker, and Stanford (MSS). This upper bound is saturated when the system under consideration is the exact holographic dual of a black hole. 
It has also been shown that an inverted harmonic oscillator (IHO) may exhibit the 
behavior of thermal energy emission, in close analogy to the Hawking radiation emitted by black holes.  
We demonstrate that the Lyapunov exponent of the IHO indeed saturates the MSS bound, 
with an effective temperature equal to the analogue black hole radiation temperature, 
%Therefore, the IHO can elegantly capture the essential features of a black hole. 
and propose using a trapped ion as a physical implementation of the IHO. 
We    
% and explore the analogue Hawking radiation through this correspondence. 
derive the corresponding out-of-time-ordered correlation function 
(OTOC) diagnosing quantum chaos, % and its Lyapunov exponent.  
%and the scattering off IHO potentials are studied 
and theoretically show, for an experimentally realizable setup, 
that the effective temperature of the trapped-ion-IHO matches the upper 
MSS bound for the speed of scrambling. 
% encapsulated by the Lyapunov exponent derived from the OTOC. }
\end{abstract}

%\pacs{85.25.Dq, 84.40.Az, 04.80.Cc, 04.62.+v, 98.80.Cq}
\baselineskip=0.45 cm
\maketitle
\newpage
%%%%%%%%%%%%%%%%%%%%%%%%%%%%%%%%%%%%%%%%%%%%%%%%%%%%%%%%%%%%
The Hawking effect \cite{Hawking1, Hawking1975}, describing particle creation by quantum vacuum fluctuations in the presence of a black hole event horizon, can be mimicked in the laboratory, as argued by Unruh \cite{PhysRevLett.46.1351}.  
%the mathematical equivalence between the wave propagation in curved spacetimes and the acoustic wave in condensed matter systems.
Since his seminal work, and with increasing experimental capabilities, a veritable explosion of 
interest occurred in what was dubbed {\em analogue gravity} \cite{analogue-gravity},    
involving, for example, Bose-Einstein condensates \cite{PhysRevLett.85.4643, BLV,PhysRevLett.91.240407, PhysRevA.69.033602,Ralf,PhysRevLett.115.025301,Carusotto_2008,BEC1, BEC2, BEC3, Renaud,PhysRevLett.124.060401,PhysRevLett.118.130404,Tyler} light in nonlinear media \cite{Philbin1367, PhysRevLett.105.203901, PhysRevLett.122.010404}, magnets \cite{PhysRevLett.118.061301,PhysRevLett.122.037203}, or superconducting circuits \cite{PhysRevLett.103.087004, PhysRevD.94.081501, PhysRevD.95.125003,PhysRevD.100.065003,superconducting-circuit}. 
%see Ref. \cite{analogue-gravity} for an extensive review.
% and a comprehensive list of references.




Gauge-gravity duality \cite{GGD1, PhysRevD.58.046004} represents a powerful %computational 
tool enabling the description of quantum gravity by certain classes of  
large-N gauge theories. 
Recently, it was conjectured that the Lyapunov exponent $\lambda_L$
characterizing the rate of growth of chaos in thermal quantum many-body systems 
%with a large number of degrees of freedom,
has a rigorous bound, $\lambda_L\leq2\pi\,T/\hbar$ \cite{bound}. 
This bound is saturated when the quantum system 
is exactly holographic dual to a black hole \cite{QBC1, Sekino_2008, Franz:2018cqi}, 
and thus may capture an essential feature of the gauge-gravity duality. 
Conversely, the bound might suggest that a %chaotic 
quantum system with given $\lambda_L$ has a minimal temperature $T\ge\hbar\lambda_L/2\pi$. 
Thus a system with $\lambda_L\neq 0$, although possessing zero temperature classically, 
%, which is classically chaotic at absolute zero %($\lambda_L\neq 0, T\rightarrow 0,\hbar\rightarrow 0$)
may exhibit thermality in the semiclassical regime \cite{PhysRevLett.122.101603}, that is
$\lambda_L\neq 0, T\neq 0$, when $\hbar\neq 0$.
%This feature is cleverly linked to the thermodynamics of black holes---a
This recalls a salient property of black holes: While being completely black classically, 
they semiclassically radiate at the Hawking temperature.  
%Studying semiclassical chaotic systems thus also 
%paves a new way for observing Hawking radiation. 


%Particle motion in an IHO potential is chaotic.  
% characteristic.}
%s very sensitive to initial conditions and hence 
%, with a certain $\lambda_L$, while 
The quantum dynamics of the IHO exhibits thermal behavior with a temperature depending
on the %classical 
Lyapunov exponent $\lambda_L$ \cite{PhysRevLett.122.101603, Betzios, Maldacena_2005}. 
The IHO model has, therefore, been used to 
explore the MSS bound on system temperature from a different angle,  
revealing its relation to analogue Hawking radiation % in the semiclassical regime
 \cite{PhysRevLett.122.101603},
and to study the scattering outside a black hole horizon quantum mechanically \cite{Betzios, Maldacena_2005}. 
However,  to experimentally verify whether the claimed duality between black holes and the fastest quantum scramblers  
indeed exists in nature, and how to extract the Lyapunov exponent within an experimentally accessible system 
are still open issues. 
To fill this gap, we propose below an implementation of the IHO with a trapped ion.
%, and demonstrate a gauge-gravity duality corresponding to the MSS bound. % in a readily accessible physical system. 
Specifically, we propose to determine the Lyapunov exponent by measuring the IHO's 
OTOC \cite{1969JETP1200L, bound}. 
%, using a new protocol which does not necessitate reversing the time evolution \cite{blocher2020measuring}. 
%Furthermore, t
To establish gauge-gravity duality, we show that IHO analogue Hawking temperature and  
Lyapunov exponent indeed fulfill $\lambda_L= 2\pi\,T/\hbar$. 

%Below, we propose an implementation of the IHO with a trapped ion, 
%and show that it realizes the gauge-gravity duality 
%in a readily accessible physical system.
%\colb{Classical one-dimensional particle motion in an inverse harmonic potential is very sensitive to initial conditions and chaotic, with a certain
%$\lambda_L$, while the corresponding quantum dynamics exhibits thermal behavior with a temperature depending 
%on $\lambda_L$ \cite{PhysRevLett.122.101603, Betzios, Maldacena_2005} .}  
%To diagnose the ion's chaotic motion, 
%we propose to measure its OTOC \cite{1969JETP1200L, bound}, 
%using a new protocol which does not necessitate reversing the time evolution. Furthermore, 
%we identify how the analogue Hawking temperature in the ion relates to the Lyapunov exponent. We demonstrate that our results are  within 
%current experimental reach. 
Trapped ions, featuring a unique level of fidelity in preparation, control, and readout of quantum states, have been extensively used to simulate relativistic quantum physics, 
such as Zitterbewegung \cite{PhysRevLett.98.253005, ZMD, PhysRevLett.120.160403}, the 
Klein paradox \cite{PhysRevLett.106.060503, PhysRevA.82.020101}, and (analogue) cosmological particle creation  
%stimulated by quantum vacuum fluctuations  
\cite{PhysRevLett.94.220401, PhysRevLett.99.201301, PhysRevLett.123.180502, Menicucci_2010}. 
We demonstrate that our proposal furnishes an implementation of exact 
gauge-gravity duality within current experimental reach for trapped ions. % control. 

The one-dimensional IHO, with Hamiltonian ($\alpha >0$)
 \begin{eqnarray}\label{Hamiltonian}
H=\frac{{p}^2}{2m}-\frac{\alpha}{2}{x}^2, 
%\quad (\alpha>0), 
\end{eqnarray}
%describes a particle with mass $m$ 
%in a one-dimensional inverted harmonic potential (with curvature $\alpha>0$) 
yields the classical trajectory 
$x(t)=c_1\,e^{\sqrt{\alpha/m}t}+c_2\,e^{-\sqrt{\alpha/m}t}$ (see discussion below for the relation of the  IHO ``particle" to the actual ion).
This deterministic evolution is exponentially sensitive to the initial condition specified by $c_1$ and $c_2$, 
and the Lyapunov exponent $\lambda_L=\sqrt{\alpha/m}$.  
%The quantum variant of the IHO satisfies the Schr\"odinger equation with the Hamiltonian,
%where $\hat{p}$ and $\hat{x}$ are the momentum and position operators, respectively. 
In the quantum version of the IHO, the OTOC 
%is a correlation of 
for momentum and position operators is given by: 
\begin{eqnarray}\label{OTOC}
C(t)=-\langle[\hat{x}(t),\hat{p}(0)]^2\rangle\sim\hbar^2e^{2\lambda_Lt}. 
\end{eqnarray}
Applying the Lyapunov exponent bound  \cite{bound}, $\lambda_L\leq2\pi\,T/\hbar$, to the 
quantum IHO predicts the existence of a lower bound on temperature 
$T\geq\hbar\lambda_L/2\pi$, as conjectured in \cite{PhysRevLett.122.101603}. 
For the classical IHO ($\hbar\rightarrow 0$), which is both nonthermal and deterministic, 
this inequality becomes an equality as $T\coloneqq 0$.  
Conversely, in the semiclassical regime, following \cite{PhysRevLett.122.101603},  
the classical Lyapunov exponent $\lambda_L$ acquires a 
quantum correction ${\cal\,O}(\hbar)$, so that the right-hand side of the inequality 
becomes $T\geq\frac{\hbar}{2\pi}[\lambda_L+{\cal\,O}(\hbar)]=\frac{\hbar}{2\pi}\lambda_L+{\cal\,O}(\hbar^2)$. 
This suggests that, at least on the semiclassical level, an effective temperature 
of ${\cal\,O}(\hbar)$ is induced, which closely resembles the situation of black hole thermodynamics: While black holes are classical solutions of general relativity, semiclassically, they create a thermal bath of radiation 
 \cite{Hawking1, Hawking1975}. 
 %they behave 
%semiclassically as if they were in a bath at the Hawking temperature \cite{Hawking1, Hawking1975}. 
%This formal similarity represents our starting point. % of our following analysis.

To establish the correspondence of black hole radiation and quantum behavior 
%chaos 
of the IHO, we consider scattering off the IHO potential. 
%\col{This S Matrix calculation is similar to the Bogoliubov calculation of Hawking or is it not?} 
Defining the light-cone-type 
operators $\hat{u}^\pm=(\hat{p}^\prime\pm\hat{x}^\prime)/\sqrt{2}$ with $\hat{p}^\prime=\hat{p}/\sqrt{m}$ and $\hat{x}^\prime=\sqrt{\alpha}\hat{x}$, we can represent 
%rewrite the Hamiltonian 
\eqref{Hamiltonian} in the $u^\pm$ basis as $H_\pm=\mp\,i\hbar\lambda_L\big(u^\pm\partial_{u^\pm}+\frac{1}{2}\big)$, respectively. 
\iffalse
\bea 
H=i\hbar\lambda_L\left(\begin{array}{cc}
%\mp\,i\hbar\lambda_L\bigg(u^\pm\partial_{u^\pm}+\frac{1}{2}\bigg).
-\bigg(u^+\partial_{u^+}+\frac{1}{2}\bigg) & 0\\
0 & i\bigg(u^-\partial_{u^-}+\frac{1}{2}\bigg)
\end{array}\right)
\ea
\fi
The in- and outgoing energy eigenstates 
are of the form $\frac{1}{\sqrt{2\pi\hbar\lambda_L}}(\pm\,u^\pm)^{\pm\,i\frac{\varepsilon}{\hbar\lambda_L}-\frac{1}{2}}\Theta(\pm\,u^\pm)$, where $\Theta(u^\pm)$ is Heaviside step function, and $\varepsilon$ is energy. 
Incoming and outgoing
states are connected by a Mellin transform, which gives the scattering matrix ($S$ matrix) 
in the form \cite{Betzios, Maldacena_2005} 
\begin{eqnarray}\label{S-Matrix}
\nonumber
S&=&\frac{1}{\sqrt{2\pi}}\exp\bigg[-i\frac{\varepsilon}{\hbar\lambda_L}\log\hbar\lambda_L\bigg]\Gamma\bigg(\frac{1}{2}-i\frac{\varepsilon}{\hbar\lambda_L}\bigg)
\\    \nonumber
&&\times
\left(
\begin{array}{cccc}
e^{-i\frac{\pi}{4}}e^{-\frac{\pi\,\varepsilon}{2\hbar\lambda_L}} & e^{i\frac{\pi}{4}}e^{\frac{\pi\,\varepsilon}{2\hbar\lambda_L}}
\\
e^{i\frac{\pi}{4}}e^{\frac{\pi\,\varepsilon}{2\hbar\lambda_L}} & e^{-i\frac{\pi}{4}}e^{-\frac{\pi\,\varepsilon}{2\hbar\lambda_L}}
\end{array}
\right), \label{Smatrix} 
\end{eqnarray}
where $\Gamma$ is the Gamma function. From the $S$ matrix, the transmission and reflection probability are respectively given by, using 
%\colp{The $S$ matrix describes the relations between the incident, reflection, and transmission waves. So 
%here T and R are the matrix elements of $S$. Their square modulus is the reflection and transmission probability.
%During the calculation, the gamma function can be simplified: 
$|\Gamma\big(\frac{1}{2}-i\frac{\varepsilon}{\hbar\lambda_L}\big)|^2=\pi\sech(\frac{\pi\varepsilon}{\hbar\lambda_L})$, 
\begin{eqnarray}\label{TR}
|{\mathcal T}(\varepsilon)|^2=\frac{1}{1+e^{2\pi\,\varepsilon/\hbar\lambda_L}},~~|{\mathcal R}(\varepsilon)|^2=\frac{1}{1+e^{-2\pi\,\varepsilon/\hbar\lambda_L}}, \label{TandR}
\end{eqnarray}
corresponding to a thermal distribution, with a temperature given by $T=\hbar\lambda_L/2\pi$. 

As illustrated in Fig. \ref{fig1}, we can understand the transmission coefficient above 
by the quantum mechanical tunneling process through a barrier 
of the particle with negative energy $\varepsilon$ moving towards the inverted harmonic potential from the left $(x\to0)$, with a tunneling probability $|\mathcal{T}(-\varepsilon)|^2$.  
%In classical mechanics, the particle moves following the classical trajectories. 
%That is to say, the negative energy particle is reflected by the potential and 
%goes back $(x\to-\infty)$, however, the positive energy particle goes through the potential toward $x\to+\infty$. If quantum corrections are considered, the negative energy particle $(\varepsilon<0)$, as a consequence of the quantum tunneling,
%can penetrate the potential, even which is forbidden for the classical mechanics case. 
%The tunneling probability is $|T(-\varepsilon)|^2$, 
%means that the ratio of the probability of taking the classical trajectory $(x\to-\infty)$ to that of %the quantum one $(x\to+\infty)$ is
%$1$ to $e^{\beta_L\varepsilon}$ with $\beta_L=1/T$. 
Similarly, the incoming positive energy particle $(\varepsilon>0)$ from the right
may be reflected by the potential quantum-mechanically, with probability $|{\mathcal R}(-\varepsilon)|^2$. 
%It means that the ratio of the probability of taking the classical 
%trajectory $(x\to+\infty)$ to that of the quantum one $(x\to-\infty)$ is $1$ to $e^{-\beta_L\varepsilon}$.

\begin{figure}[t]
\centering
\includegraphics[width=0.45\textwidth]{BH}
\caption{(Color online) (a) Scattering off an effective potential approximately equal to an inverse harmonic potential, outside a black hole, where $r_\ast$ is the ``tortoise coordinate", and the event horizon is located at $r_\ast\rightarrow-\infty$. 
(b) Scattering via the inverse harmonic potential leads to classical trajectories of incoming particles  (solid lines) and particle tunneling (broken lines) near the hyperbolic fixed point $(x, p)=(0, 0)$ in phase space. The dotted lines are the separatrices $(\varepsilon=0)$, and $\beta=1/T$. }
% is the inverse of temperature.}
\label{fig1}
\end{figure}

A relativistic particle, when close to a black hole horizon (pulled by external  
scalar or electromagnetic forces so that it does not fall into the black hole),  
possesses an effective Lagrangian corresponding to the 
Hamiltonian \eqref{Hamiltonian} \cite{PhysRevD.95.024007}. 
%\colb{A relativistic particle moving close to a  black hole horizon 
%which is pulled by scalar or electromagnetic forces toward outside the horizon 
%possesses an effective Lagrangian corresponding to the nonrelativistic Hamiltonian \eqref{Hamiltonian} 
%\cite{PhysRevD.95.024007}. }
The maximal Lyapunov exponent of the particle's motion is  
found to be the surface gravity of the black hole when the total potential presents unstable maximum \cite{PhysRevD.95.024007}. The Lyapunov exponent is then given by 
$\lambda_L=\kappa=\frac{1}{2}\sqrt{g^\prime(r_0)f^\prime(r_0)}$, for the metric $ds^2=-f(r)dt^2+dr^2/g(r)+r^2(d\theta^2+\sin^2\theta\,d\phi^2)$. Here $g^\prime(r_0)$ and $f^\prime(r_0)$ are derivatives with respect to $r$ evaluated at the horizon $r=r_0$. 
%In classical sense, black hole is completely black,
%and no radiation from the black hole can be expected. While in quantum sense, it behaves as a thermal bath with the temperature equal
%to its surface gravity. 

Numerous variants of (re-)deriving Hawking's result exist  
\cite{Hawking1, Hawking1975, Gibbons, Parikh, Robinson, Satoshi, Peet:2000hn, STROMINGER199699, PhysRevD.77.024031}, ranging, e.g., from Hawking calculating the Bogoliubov coefficients between the quantum scalar field modes of the ``in" and ``out" vacuum states \cite{Hawking1, Hawking1975}, to an open quantum system approach \cite{PhysRevD.77.024031}. Here, we 
follow Parikh and Wilzcek \cite{Parikh}, employing the tunnelling probability through a classically forbidden region.
% to describe Hawking radiation. 

Consider our particle tunneling through an effective black hole potential due to a general spherically symmetric  metric $ds^2=-f(r)dt^2+dr^2/g(r)+r^2(d\theta^2+\sin^2\theta\,d\phi^2)$. 
On the semiclassical level under consideration here, we can employ the 
WKB approximation to calculate the tunnelling (transmission) probability through the horizon. 
To leading order in $\hbar$, it is given by \cite{Kerner2008}
\begin{eqnarray}
\Gamma=\exp\bigg[-\frac{4\pi}{\sqrt{g^\prime(r_0)f^\prime(r_0)}}\varepsilon\bigg]\coloneqq\exp[-\varepsilon/T]. 
\end{eqnarray}
%This means that  particles can tunnelling, which is forbidden in classical mechanics, from inside to outside the horizon in semi-classical mechanics, understood as the intriguing 
The above tunneling rate implies Hawking radiation with 
temperature $T=\sqrt{g^\prime(r_0)f^\prime(r_0)}/4\pi$.
We conclude that IHO and black hole thus share the 
property that while classically not behaving as if being in thermal equilibrium with a bath, thermality is acquired when 
quantum mechanics is taken into account.  
This formal analogy forms the basis of our proposal to observe analogue 
Hawking radiation in an experimentally feasible  quantum IHO system, to which we now proceed.


We propose using a trapped ion to experimentally implement the IHO.
In our scenario, we monitor one of the oscillation directions of the ion,  
%\colp{This is commonly done for trapped ion experiment. Through adjusting the trap potential, one can confine the ion's spatial motion. 
%For almost all the trapped ion experiments, ion harmonically oscillates in one-direction, while the motion in the other two directions 
%is confined.} 
with associated momentum and position operators in the lab frame denoted as $\hat{P}$ and $\hat{X}$, respectively.
In particular, the ion's oscillation frequency $\omega_0$ can be 
periodically modulated by an applied voltage on the trap electrodes \cite{RevModPhys.75.281}.
The corresponding Hamiltonian is given by $H=\frac{\hat{P}^2}{2M}+\frac{1}{2}M\omega^2(t)\hat{X}^2$,
where $\omega(t)=\omega_0[1-\xi\cos(\omega_mt+\phi)]$ denotes the time-dependent frequency, $\xi$ is the modulation depth, $\phi$ is the initial phase of the modulation, and $M$ is the mass of the ion. 
Imposing $0<\xi\ll1$, and $\omega_m=2\omega_0$, 
the Hamiltonian, in rotating-wave approximation, 
%$\omega_0\gg1$, 
%\col{omega is not dimensionless,so what is ``1" here!!} 
%\colp{A mistake here, during the calculation, we obtained $H\approx-\frac{1}{4}\hbar\omega_0\xi[a^2e^{i\phi}+a^{\dagger2}e^{-i\phi}+a^2e^{-i(4\omega_0t+\phi)}+a^{\dagger2}e^{i(4\omega_0t+\phi)}]$, the 
%rotating-wave approximation condition requires $4\omega_0\gg\frac{1}{4}\omega_0\xi$, simplified to $\xi\ll16$, which is 
%satisfied by the $\xi$ assumption previous. So here I think we have no need to write its detailed form.}
can be rewritten in the form of the IHO 
in Eq.~\eqref{Hamiltonian}, with $m=2M/\xi$, and $\alpha=\frac{1}{4}m\omega_0^2\xi^2=m\omega^2$. The momentum and position 
operators have been redefined by $\hat{p}=-i\sqrt{\frac{m\omega\hbar}{2}}(\hat{a}e^{i\phi/2}-\hat{a}^\dagger\,e^{-i\phi/2})$,
and $\hat{x}=\sqrt{\frac{\hbar}{2m\omega}}(\hat{a}e^{i\phi/2}+\hat{a}^\dagger\,e^{-i\phi/2})$, respectively. Note that the phase $\phi$ is tunable by appropriately choosing the initial phase of driving. 



\iffalse
and that the rotating-wave approximation is 
satisfied, the Hamiltonian \eqref{H-ion} then is equivalent 
to the Hamiltonian \eqref{Hamiltonian}, but with a scaling factor $\xi$ relating to 
the driving strength, i.e.,
$H_\text{rot}=\frac{\xi}{2}\big[\frac{\hat{p}^2}{2m}-\frac{1}{2}m\omega^2\hat{x}^2\big]
=\frac{\hat{p}^2}{2m^\prime}-\frac{\alpha^\prime}{2}\hat{x}^2$, 
where $m^\prime=2m/\xi$, and $\alpha^\prime=\frac{1}{4}m^\prime\omega^2\xi^2=m^\prime\omega^{\prime2}$. Therefore, the 
IHO Hamiltonian in Eq. \eqref{Hamiltonian} could be implemented in a 
trapped ion setup. 
We omit for simplicity the primes on various quantities in what follows. 
%For simplification,we still use $m$, $\omega$, and $\alpha$ instead of $m^\prime$, $\omega^%\prime$, and $\alpha^\prime$ hereafter.
\fi


To diagnose the \colb{quantum} chaotic motion of the IHO, we propose to measure its OTOC using the 
protocol proposed in Ref.~\cite{blocher2020measuring} 
which does not require the reversal of time evolution. Here, 
the OTOC is defined by  $C(t)=\mathrm{Tr}[\rho_0W^\dagger(t)V^\dagger\,W(t)V]$, 
where $W(t)=U^\dagger(t, 0)WU(t, 0)$ and $V=V(0)$ are operators evaluated in the
Heisenberg picture at times $0$ and $t$, respectively \cite{1969JETP1200L, bound}. 

Assuming the operator $V$
to be a projection operator onto an initially pure state, $V=\rho_0$  \cite{blocher2020measuring},  the OTOC can be reduced to $C(t)=|\langle\,W(t)\rangle|^2$.
To measure the OTOC, we therefore need 
to prepare a pure initial state for the degree of freedom of the ion's motion. For our 
experimental proposal, we use for concreteness a single $^9\mathrm{Be}^+$ ion trapped in a strong radio-frequency Paul trap with a pseudopotential trap frequency 
of $\omega_0/2\pi=10\,\mathrm{MHz}$. The motional ground state can be prepared 
by a two-stage laser cooling process. By Doppler cooling, all the motional modes of the ion could be cooled down to near the Doppler limit (with an average phonon number $\bar{n}$ of typically a few to ten quanta), through driving the 
$^2S_{1/2}$ to $^2P_{3/2}$ dipole transition. %Subsequently, w
We can further cool the motional mode of interest to its ground state using sideband cooling with stimulated Raman transitions between the ion's two hyperfine ground states, i.e., between $^2S_{1/2}(F=2, m_F=2)$ and $^2S_{1/2}(F=1, m_F=1)$, denoted by $|\downarrow\rangle$ and $|\uparrow\rangle$, respectively, and are separated by $\sim1.25\,\mathrm{GHz}$ \cite{PhysRevA.42.2977, PhysRevLett.75.4011}. 
\iffalse
%This can be realized using sideband cooling with stimulated Raman transitions \cite{PhysRevLett.75.4011} between the ion's two 
%hyperfine ground states, e.g., between $^2S_{1/2}(F=2, m_F=2)$ and $^2S_{1/2}(F=1, m_F=1)$ of the $^9\mathrm{Be}^+$ ion, which
%are denoted by $|\downarrow\rangle$ and $|\uparrow\rangle$, respectively. 
\fi
The spin initialization and ground state cooling allow us to prepare the ion in the 
$|S=\downarrow, n=0\rangle$ state eventually. By applying appropriate laser pulse sequences on blue and red sidebands or the carrier, 
%frequency, 
we can create  arbitrary quantum superpositions of Fock states \cite{PhysRevA.55.1683, PhysRevA.57.2096, PhysRevLett.90.037902, PhysRevLett.92.113004}. 
%By simply applying a sequence of Rabi $\pi$ pulses of laser radiation on the blue sideband, red sideband or the carrier, we can create higher
%$n$ Fock states, denoted by $\rho_0=|n\rangle\langle\,n|$ \cite{PhysRevLett.76.1796, Schneider_2012}. 

For the proposed experiment, we assume that the perturbation operator $W$ is an ion displacement measurement, 
$\hat{x}=\sqrt{\frac{\hbar}{2m\omega}}(\hat{a}+\hat{a}^\dagger)$, %\cite{RevModPhys.84.621}, 
and
the initial pure state is prepared as, $| \Psi_\text{in}\rangle=\frac{1}{\sqrt{2}}(|n\rangle+|n+1\rangle)$, \cite{PhysRevLett.90.037902, PhysRevLett.92.113004}.
%i.e., \colp{$W=\hat{a}^\dagger\,\hat{a}$}, with \colp{$\hat{a}$ and $\hat{a}^\dagger$} being the annihilation and creation operators of phonons, %respectively. 
%decoherence can be modeled as a perturbation operation $W$, 
%which is equivalent to a projector on the 
%Fock state, $W=a^\dagger\,a$, \colb{and amounts to a dephasing jump operator}, 
%with $a$ and $a^\dagger$ being the annihilation and creation operators of phonons, respectively.
We then readily find that the OTOC is given by 
\begin{eqnarray}\label{OTOCn}
\nonumber
C(t)&=&\left|\langle \Psi_\text{in}|\,U^\dagger(t, 0)\hat{x}U(t, 0)|\Psi_\text{in}\rangle\right|^2
\\&=&
\frac{(n+1)\hbar}{2m\omega}\cosh^2(\lambda_Lt),
\end{eqnarray}
where $U(t, 0)=\exp[-iHt/\hbar]$. Measuring  OTOC then reduces to determining  directly 
accessible ion-motional quadratures  
\cite{ZMD, PhysRevLett.104.100503, PhysRevA.100.062111}. 
We find that
the OTOC exhibits exponential growth, $\propto e^{2\lambda_Lt}+(\text{terms growing more slowly than}\,e^{2\lambda_Lt})$, with a Lyapunov exponent 
%\col{??} 
$\lambda_L=\sqrt{\alpha/m}$. 
%\col{??}
Fig.~\ref{fig2} shows the OTOC for various pure initial states and Lyapunov exponents.  

\begin{figure}[t]
\centering
\subfigure[]{\includegraphics[width=0.2\textwidth]{F1}}\hspace*{2em}
\subfigure[]{\includegraphics[width=0.2\textwidth]{F2}}
\subfigure[]{\includegraphics[width=0.2\textwidth]{F4}}\hspace*{2em}
\subfigure[]{\includegraphics[width=0.2\textwidth]{F5}}
\subfigure[]{\includegraphics[width=0.25\textwidth]{F3}}
\caption{(Color online) (a) The OTOC in Eq. \eqref{OTOCn} for various pure initial Fock superposition states with driving strength $\xi=0.01$, and $x_0=\sqrt{\hbar/2m\omega}$ is vacuum position uncertainty; (b) the OTOC in Eq.~\eqref{OTOCn} for different Lyapunov exponent parameters (driving strengths) $\xi$; the initial pure state is $|\Psi_\text{in}\rangle=(|0\rangle+|1\rangle)/\sqrt{2}$; (c) the transmission and reflection probability 
in Eq.~\eqref{TR-pro} as function of ion displacement $x$, with fixed width $\Delta=\hbar\omega$, $\xi=0.01$, initial energy 
$\varepsilon_0=-0.5\hbar\omega$, and at time $t=2.86\,\mathrm{\mu\,s}$; (d)  
%the transmission and reflection probability in Eq. \eqref{TR-pro} 
transmission and reflection probability as a function of time $t$, with 
%with identical parameters $\Delta$, $\xi=0.01$, and $\varepsilon_0=-0.5\hbar\omega$. 
ion displacement $x=1.35\,\mathrm{nm}$. 
(e) the fidelity of the prepared initial state \eqref{incident-wave} for various truncation 
levels $L$ of the Fock basis.
Everywhere, the oscillator frequency is taken to be 
$\omega_0/2\pi=10\,\mathrm{MHz}$.} \label{fig2}
\end{figure}



To observe the Hawking radiation associated to the S matrix \eqref{Smatrix}, 
we need to prepare a suitable incident wave packet. 
Specifically, we prepare the initial state of the ion's motion as a superposition of energy states, 
$|\Psi_\text{in}\rangle=\int\,d\varepsilon\,f(\varepsilon)|\varepsilon\rangle$.
%\colp{In order to prepare this initial state in experiment, we calculate the representation of this initial state in the Fock basis, see below. Since it is %convenient for us to prepare different Fock states and their superposition by the blue sideband, red sideband and the carrier in experiment. Therefore, if we know the representation of this initial sate in the Fock basis, then we can use different combination (i.e., different pulses) of the blue sideband, red sideband and the carrier to prepare the corresponding Fock state of the initial state. I will give the illustration in the following, including the Refs. [63, 64].}
When the distribution $f(\varepsilon)$ is assumed to be a Gaussian centered at an energy $\varepsilon_0$ with small width $\Delta$, i.e., 
$f(\varepsilon)=\frac{1}{\sqrt{2\pi^{3/2}\Delta}}\exp\big[-\frac{(\varepsilon-\varepsilon_0)^2}{2\Delta^2}\big]$, the incident wave packet can be written as \cite{BARTON1986322}
\begin{eqnarray}\label{incident-wave}
\Psi_\text{in}(x,t)=i{|x|^{-1/2}}e^{-i\varepsilon_0t-i\Phi_0}
F\big(t+\log[\sqrt{2}|x|]\big),
\end{eqnarray}
where $F(z)=(\Delta/\pi^{1/2})^{1/2}\exp\{-z^2\Delta^2/2\}$, $\Phi_0=\varepsilon_0\log\sqrt{2}|x|+x^2/2+\varphi_0/2+\pi/4$, 
with $\varphi_0=\arg\Gamma(\frac{1}{2}-i\varepsilon_0)$. 
Generating arbitrary harmonic-oscillator states has been experimentally demonstrated recently \cite{PhysRevLett.90.037902}. 
To prepare this wave packet in experiment, it is convenient 
%\colp{Since people usually control and manipulate the trapped ion system by laser to implement the blue sideband, red sideband and the carrier. These operations allow us to excite and deexcite the phonon Fock state, and simultaneously flip the ion's internal state. Thus if we know the form of the initial state in the Fock basis, it is more convenient for us to prepare it with appropriate combination of the blue sideband, red sideband and the carrier.}
to represent it in a phononic Fock basis, 
$|\Psi_\text{in}\rangle=\sum^\infty_{n=0}\langle\,n|\Psi_\text{in}\rangle|n\rangle$. 
We note that the overlap $\langle\,n|\Psi_\text{in}\rangle$
is nonzero only when $n$ is even. Using the experimental techniques of Refs.~\cite{PhysRevLett.90.037902, PhysRevLett.92.113004},  
the incident wave packet \eqref{incident-wave} can in principle  
be prepared experimentally. % in the Fock space.
In Fig.~\ref{fig2}, we plot the fidelity 
$F_\text{in}=\sum^L_{n=0}|\langle\,n|\Psi_\text{in}\rangle|^2$ for preparing the incident wave packet, with $L$ being the truncation 
level in Fock space. 
%of this incident wave packet. % by different truncation in the Fock space. 
For the initial state \eqref{incident-wave}, we derive an approximate analytical  expression 
for the transmitted and reflected probability densities as follows \cite{BARTON1986322}:
\begin{eqnarray}\label{TR-pro}
\nonumber
|\Psi_R|^2&=&|{\mathcal R}(\varepsilon_0)|^2\frac{1}{|x|}\left|F\left(t-\log[\sqrt{2}|x|]-\varphi^\prime(\varepsilon_0)\right)\right|^2,
\\ 
|\Psi_T|^2&=&|{\mathcal T}(\varepsilon_0)|^2\frac{1}{x}\left|F\left(t-\log[\sqrt{2}x]-\varphi^\prime(\varepsilon_0)\right)\right|^2. 
\end{eqnarray}
%where $F(z)=(\Delta/\pi^{1/2})^{1/2}\exp\{-z^2\Delta^2/2\}$, and 
Here, $\varphi^\prime =d\varphi/d\varepsilon$, with  %denotes the derivative of 
$\varphi(\varepsilon)=\arg\Gamma(\frac{1}{2}-i\varepsilon)$.  
%with respect to its argument. 
%For simplicity, w
We have taken the units for mass, length, time, and energy as $2m$, 
$(\hbar/2m\omega)^{1/2}$, $\omega^{-1}$, and $\hbar\omega$, respectively. 
Integrating Eq.~\eqref{TR-pro} over space, we obtain the energy-dependent transmission and reflection probabilities in Eq.~\eqref{TR}.


In an experiment, the quantum mechanical ion motion can be detected by observing the evolution of its internal level populations  
according to the Hamiltonian % interaction 
$H_\mathrm{I}=\hbar\Omega\sigma_y(\hat{a}e^{i\phi/2}+\hat{a}^\dagger\,e^{-i\phi/2})$  \cite{PhysRevLett.76.1796}. 
This interaction can be implemented, in the Lamb-Dicke limit, 
by driving the ion with both blue and red sidebands, $H_b=\hbar\eta\Omega_b(\hat{a}^\dagger\sigma^+e^{i\phi_b}+\hat{a}\sigma^-e^{-i\phi_b})$ and 
$H_r=\hbar\eta\Omega_r(\hat{a}\sigma^+e^{i\phi_r}+\hat{a}^\dagger\sigma^-e^{-i\phi_r})$, respectively, 
%simultaneously on the ion, i.e., 
and thus with the total Hamiltonian  $H_b+H_r
=\hbar\Omega(\hat{a}\,e^{i\phi_-}+\hat{a}^\dagger\,e^{-i\phi_-})(\sigma^+e^{i\phi_+}+\sigma^-e^{-i\phi_+})$. Here,  
we set $\eta\Omega_b=\eta\Omega_r=\Omega$, $2\phi_\pm=\phi_r\pm\phi_b$ by tuning the amplitude and phase of the applied driving field for the sidebands, and $\sigma^\pm=(\sigma_x\pm\,i\sigma_y)/2$.
%Here $\sigma_y$ is the Pauli operator that acts on the two 
%internal levels of the ion, and $\Omega$ is the Raman coupling frequency. 
To measure the scattering state, we first need to prepare the internal levels of the ion as  
$|+\rangle=1/\sqrt{2}(|\uparrow\rangle+|\downarrow\rangle)$,
and then apply the Raman-induced interaction coupling $H_\mathrm{I}$. 
We then record the probability $P_\downarrow(t)$ of occupation in $|\downarrow\rangle$.
The expected signal is $P_\downarrow(t)=\int\big(1-\sin(2\Omega\,tx)\big)|\Psi(x)|^2dx$ and 
the scattering probability $|\Psi(x)|^2$ in Eq. \eqref{TR-pro}
is then obtained by applying an inverse Fourier transform to $P_\downarrow(t)$. 




Taking the initial energy, $\varepsilon_0<0$, 
we plot transmission and reflection probabilities in Fig.~\ref{fig2}. 
Negative $\epsilon_0$ implies that the particle
%is smaller than the inverse harmonic potential. From the perspective
classically does not penetrate the barrier when coming from the left; 
%it will be scattered to the right and back to the left, and has no possibility to penetrate the %potential. 
Fig.~\ref{fig2} however clearly shows that both nonzero transmission and reflection probability 
exist simultaneously. 
%It suggests that 
%quite unlike the classical case, due to the quantum tunneling, the particle can penetrate the potential, although that is bigger than
%the particle's kinetic energy, and arrive the right region that would not be allowed by the laws of classical physics. Interestingly,
The transmission probability, after integrating over position, 
satisfies a thermal distribution with temperature $T=\hbar\lambda_L/2\pi$. 
This mathematical relation between the temperature $T$ and the Lyapunov exponent $\lambda_L$ suggests that the conjectured MSS bound \cite{bound} is indeed satisfied for the ion's motion. 
We have thus shown that the ion trap IHO is the dual of a black hole. 
%Therefore, the relevant physics of the IHO discussed above can be understood as the %analogue Hawking radiation.


Rewriting Eq.~\eqref{Hamiltonian} with phononic annihilation and creation operators, 
we find $H=-\frac{1}{2}\hbar\lambda_L(\hat{a}^{\dagger2}+\hat{a}^2)$, which  
corresponds to a squeezing operation \cite{Loock}. %,WANG20071}. %,RevModPhys.84.621}.   
%Therefore, the harmonic oscillator wave function remains an even Gaussian function.
We can thus alternatively interpret the analogue Hawking radiation via a squeezed state of the ion's motion. 
Phonon number increases as  
$\langle\hat{n}\rangle=\sinh^2(\lambda_Lt)$, and the squeezed vacuum state population distribution 
%for a vacuum squeezed state 
is restricted to even states, $P_{2n}=(2n)!\tanh^{2n}(\lambda_Lt)/(2^nn!)^2\cosh(\lambda_Lt)$  \cite{WANG20071}. 
The squeezed state can be detected by the evolution of the ion's internal levels under a Jaynes-Cummings interaction \cite{PhysRevLett.76.1796}. %Taking the squeezing parameter \colp{$\lambda_Lt\approx0.88$}, 
%corresponds to a squeezed state with $\langle\hat{n}\rangle\approx1$. 
The population distribution is shown in Fig. \ref{fig3}, and 
compared with a slightly thermal state close to the vacuum.  
The increased population of higher Fock states   
%\colp{This distribution is for the squeezed state. ``increased" relative to the initial vacuum state. Since this state is obtained 
%by applying the Hamiltonian operation (1) on the initial vacuum, at the beginning there is only vacuum, no other Fock states. However, %after the Hamiltonian action, other nonvacuum Fock states appear with certain probabilities.}
are evidence of squeezing; only even states populated corresponds to the creation of pairs of phonons. 


\begin{figure}[b]
\centering
\includegraphics[width=0.33\textwidth]{F6}
\caption{(Color online) Phonon number distribution $P_{n}$ for a squeezed state (yellow)
and thermal state (light blue) at very low temperatures (average phonon number $\bar{n}=0.02$). %, approaching the vacuum state}. 
The squeezing parameter is assumed to be $\lambda_Lt\approx0.88$. 
The vacuum squeezed state is created by the action of the Hamiltonian \eqref{Hamiltonian} on
the phonon vacuum. 
Although there are no phonons at the beginning, phonons are created during the evolution, as a result of analogue Hawking 
radiation.\label{fig3} }
\end{figure}

  

%%%%%%%%%%%%%%%%%%%%%%%%%%%%%%%%%%%%%%%%%%%%%%%%%%%%%%%%%%%%%%%%%%%%%%

In conclusion, we propose using well established quantum optical tools 
for trapped ions to implement the IHO as the dual of a black hole. 
Chaotic quantum motion of the ion can be accessed by measuring the OTOC, 
which exhibits to leading order exponential growth, $\propto e^{2\lambda_Lt}$,  
with Lyapunov exponent $\lambda_L=\sqrt{\alpha/m}$, 
and analogue Hawking radiation temperature $T=\hbar\lambda_L/2\pi$. 
%We observe this intriguing phenomenon as a concept of analogue Hawking radiation.
% can transfer insights between the astronomical and the mesoscopic realms. 

Numerous theoretical  approaches employed 
the IHO concept to establish links between quantum mechanics and general relativity see, e.g.,  
\cite{PhysRevLett.122.101603, PhysRevLett.123.156802, Betzios, Maldacena_2005, Morita:2018sen, PhysRevD.101.026021}.
Our IHO implementation proposal paves the way to experimentally 
explore these theoretical concepts in a highly controllable quantum optical environment.  
%Finally, %from a different angle, 
%the proposed IHO system potentially 
%captures signature features of both saddle-potential scattering of quantum Hall 
%quasiparticles and radiation phenomena in gravity \cite{PhysRevLett.123.156802}.  
%It thus provides a unique perspective to explore the connections of black hole and quantum %Hall physics.
%Our proposal is therefore expected to study the relevant physics, both gravitational and %quantum Hall physics, even some deep essential physics that links the two realms, in %experiment.

%%%%%%%%%%%%%%%%%%%%%%%%%%%%%%%%%%%%%%%%%%%%%%%%%%%%%%%%%%%%%%%%%%%%%%%%%%%%%%%%%%%%%%%%%

\begin{acknowledgments}
This work was supported by the National Key R\&D Program of China (Grant No. 2018YFA0306600), the CAS (Grants No. GJJSTD20170001 and No. QYZDY-SSW-SLH004), and Anhui Initiative in Quantum Information Technologies (Grant No. AHY050000).  
ZT was supported by the National Natural Science Foundation of 
China under Grant No. 11905218, and the CAS Key Laboratory for Research in Galaxies and Cosmology, Chinese Academy of Science (No. 18010203). YL. acknowledges support from the National Natural Science Foundation of China (Grant No. 11974330) and the Fundamental Research Funds for the Central Universities. 
URF has been supported by the NRF of Korea under Grant No.~2020R1A2C2008103. 

\end{acknowledgments}

%\newpage

\bibliography{ahr12}

\end{document}
